\section{A coefficient height register removes the common multiple block}

The previous commutative source assigned label $j$ the factor $1+jM$,
then required $M$ to be divisible by every positive integer up to a label
bound. That made the evaluated factors pairwise coprime, but introduced a
second bounded-universal certificate. The following alternative dispenses
with pairwise coprimality entirely. Its arithmetic ingredient is an elementary
bounded-height root argument; the proposed improvement is its application to
this proof source, not a claim of historical priority for that argument.

\paragraph{Polynomial identity from one integer evaluation.}
For a finite list of padded rules $(p_i,q_i,c_i)$, $i<N$, put
\[
 A(X)=\prod_{i<N}(1+c_iX),\qquad
 B(X)=(1+X)\prod_{i<N}(1+p_iX)(1+q_iX),
\]
\[
 R=2\prod_{i<N}\bigl((1+p_i)(1+q_i)+c_i\bigr).
\]
All labels are nonnegative. If $M>R$, then
\[
 A(M)=B(M)\quad\Longleftrightarrow\quad A=B\text{ in }\mathbb Z[X].
\]
Indeed, each factor in the definition of $R/2$ is at least both
$1+c_i$ and $(1+p_i)(1+q_i)$. Thus $A(1)\le R/2$ and $B(1)\le R$.
Their coefficients are nonnegative, so every coefficient of $A-B$ has
absolute value at most $R<M$. If $a_k$ were its first nonzero coefficient,
$(A-B)(M)=0$, divided by $M^k$ and reduced modulo $M$, would give
$M\mid a_k$, a contradiction.

Distinct positive labels $j$ give distinct rational roots $-1/j$ of the
linear factors. The polynomial identity therefore says exactly
\[
 \#\{i:c_i=j\}=\mathbf1_{j=1}+
 \#\{i:p_i=j\}+\#\{i:q_i=j\}\qquad(j>0).
\]
Zero labels contribute the unit factor. Consequently the same
single-conclusion Horn-flow theorem recovers derivability of target $1$.
For the calibrated numerical C222 calculus this is equivalent in ZFC to a
ZFC refutation. There is no separate primality or coprimality hypothesis.

\paragraph{Two registers and no label bounds.}
Use natural registers $S_i,R_i$ and require
\[
 S_0=M+1,\quad S_N=1,\quad R_0=2,\quad M=R_N+1,
\]
\begin{align*}
 S_{i+1}(1+c_iM)&=S_i(1+p_iM)(1+q_iM),\\
 R_{i+1}&=R_i\bigl((1+p_i)(1+q_i)+c_i\bigr).
\end{align*}
The first recurrence preserves strict positivity from $S_0>0$, so
multiplying and cancelling yields $A(M)=B(M)$. The second recurrence
provides the required coefficient bound exactly. This proves soundness
without any chronological semantics for the encoded list.

For completeness, fix a finite proof tree and list its nodes by expanding
open obligations. Compute the height product first, set $M=R_N+1$, and
then let $S_i$ be the product of $1+jM$ over its open obligations. Both
recurrences and all endpoints follow. Choosing $M$ after the finite rule
list but before the $S_i$ avoids a circular bound on a code already
containing $M$.

\paragraph{The literal local polynomial.}
Let $H(p,q,c;\mathbf w)$ be the inherited nonnegative degree-$50$ Horn
polynomial, with seventeen local witnesses. With $M=M_0+1$, define
\begin{align*}
 E&=S'(1+cM)-S(1+pM)(1+qM),\\
 F&=R'-R\bigl((1+p)(1+q)+c\bigr),\\
 L&=H+E^2+F^2.
\end{align*}
This uses twenty local natural witnesses, namely $p,q,c,\mathbf w$, and
five exposed coordinates $S,S',R,R',M_0$. Its exact degree remains $50$:
$E$ has degree five and $F$ degree three, so the additional squares cannot
cancel the degree-$50$ homogeneous part of $H$. The JSON supplies every
Horn factor and both new residuals literally.

\paragraph{One bounded block.}
Write $\beta(x,y,i)=x\bmod(1+(i+1)y)$. The five outer coordinates are
$N,x,z,y,M_0$. The endpoint conditions are
\[
 \beta(x,y,0)=M_0+2,\quad \beta(x,y,N)=1,\qquad
 \beta(z,y,0)=2,\quad \beta(z,y,N)=M_0.
\]
The only bounded-universal block is
\[
 \forall i<N\;\exists\mathbf u\in\mathbb N^{20}:\quad
 L\bigl(\beta(x,y,i),\beta(x,y,i+1),
         \beta(z,y,i),\beta(z,y,i+1),M_0;\mathbf u\bigr)=0.
\]
Both finite register sequences admit a common $y$: take $y$ larger than
all their entries and divisible by $N!$, and apply the Chinese remainder
theorem separately to $x$ and $z$. This divisibility is used only to prove
existence of the beta encoding. It is not an extra input condition and
it is not a condition on $M$.

Compared with the preceding source, the label-bound parameter, three
local slack witnesses, and the entire common-multiple bounded block are
removed. A second accumulator stream is introduced. The number of outer
parameters stays five. The supplied closed Lean constructor uses the
current remainder, bounded-forall, and quartic APIs, but has not been
elaborated or expanded. These are source-level counts, \emph{not} a
claim that the completed consistency equation has twenty variables or
that its compiled size beats the existing frontier.

\paragraph{Checks and limitations.}
The independent finite checks include $24\,844$ height fixtures, $100$
constructed proof-tree paths, twelve simultaneous beta encodings, $800$
literal residual comparisons, and twelve exact inherited Horn
sum-of-squares identities. At $M=1$, the rule with premise $1$ and
conclusion $3$ supplies the spurious unchecked accumulator path $2\to1$;
the new height bound rejects it because $R_N=10$. This negative control
shows why discarding the common-multiple block without replacing its
collision protection would be unsound. The global ordinary polynomial
remains unmaterialized, and no new proof-assistant certificate is claimed.

\paragraph{A degree-oriented selector variant.}
There is a further local reduction without increasing the old source's
local witness count. The twelve inherited branch factors have degrees
$2,4,4,4,4,4,4,4,4,8,4,4$. Partition them into six groups of degree
$8$, the truth group of degree $2$, and one unused group assigned value
$1$. Let $G_0,\ldots,G_7$ be these nonnegative products. Introduce three
selectors $b_0,b_1,b_2$ and their multilinear Boolean indicators
\[
 \chi_i(\mathbf b)=\prod_{k=0}^2
 \begin{cases}b_k,&\text{bit }k\text{ of }i\text{ is }1,\\
 1-b_k,&\text{otherwise}.
 \end{cases}
\]
Replace the Horn product by
\[
 \sum_{k=0}^2[b_k(b_k-1)]^2+
 \sum_{i=0}^7G_i\chi_i(\mathbf b)^2.
\]
All summands are nonnegative. A zero forces Boolean selectors and hence
one vanishing selected group; the unused group cannot be selected.
Conversely, any vanishing original branch admits a selector assignment.
Its exact degree is $8+6=14$. The leading terms cannot cancel because
they are the common monomial $b_0^2b_1^2b_2^2$ times a sum of nonzero
nonnegative leading forms of products of sums of squares. Adding $E^2$
and $F^2$ preserves this degree.

Thus the degree-oriented local source has twenty-three witnesses,
matching the old source, but degree fourteen instead of fifty and only
one bounded block instead of two. A separate one-selector tradeoff uses
twenty-one witnesses at degree twenty. Both have literal JSON and
uncompiled constructor files. Their finite selector audits cover
$34\,768$ and $45\,556$ assignments respectively; all three generated
constructors also pass bracket and variable-range checks, which are not
Lean elaboration.
